% Part: normal-modal-logic
% Chapter: syntax-and-semantics
% Section: language-modal-logic

\documentclass[../../../include/open-logic-section]{subfiles}

\begin{document}

\olfileid{nml}{syn}{lan}

\olsection{Basa Logika Modal Dhasar}

\begin{defn}
Basa dhasar logika modal ngemot
\begin{enumerate}
  \tagitem{prvFalse}{Konstanta proposisional kanggo !!{falsity}~$\lfalse$.}{}
  \tagitem{prvTrue}{Konstanta proposisional kanggo !!{truth}~$\ltrue$.}{}
  \item Himpunan kang !!{denumerable}s lan ngemot !!{propositional variable}s: $\Obj
    p_0$, $\Obj p_1$, $\Obj p_2$, \dots
  \item Konektif proposisional: \startycommalist
  \iftag{prvNot}{\ycomma $\lnot$ (negasi)}{}%
  \iftag{prvAnd}{\ycomma $\land$ (konjungsi)}{}%
  \iftag{prvOr}{\ycomma $\lor$ (disjungsi)}{}%
  \iftag{prvIf}{\ycomma $\lif$ (!!{conditional})}{}%
  \iftag{prvIff}{\ycomma $\liff$ (!!{biconditional})}{}.
  \tagitem{prvBox}{Operator modal $\Box$.}{}
  \tagitem{prvDiamond}{Operator modal $\Diamond$.}{}
\end{enumerate}
\end{defn}

\begin{defn}
\emph{!!^{formula}s} ing basa modal dhasar ditetepake
  kanthi induktif kaya mangkene:
\begin{enumerate}
\tagitem{prvFalse}{$\lfalse$ iku !!{formula} atomik.}{}

\tagitem{prvTrue}{$\ltrue$ iku !!{formula} atomik.}{}

\item Saben !!{propositional variable} $\Obj p_i$ iku !!{formula} (atomik).

\tagitem{prvNot}{Yen $!A$ iku !!a{formula}, mula $\lnot !A$ iku
  !!a{formula}.}{}

\tagitem{prvAnd}{Yen $!A$ lan $!B$ iku !!{formula}s, mula $(!A \land
  !B)$ iku !!a{formula}.}{}

\tagitem{prvOr}{Yen $!A$ lan $!B$ iku !!{formula}s, mula $(!A \lor !B)$
  iku !!a{formula}.}{}

\tagitem{prvIf}{Yen $!A$ lan $!B$ iku !!{formula}s, mula $(!A \lif !B)$
  iku !!a{formula}.}{}

\tagitem{prvIff}{Yen $!A$ lan $!B$ iku !!{formula}s, mula $(!A \liff !B)$
  iku !!a{formula}.}{}

\tagitem{prvBox}{Yen $!A$ iku !!a{formula}, mula $\Box !A$ iku
  !!a{formula}.}{}

\tagitem{prvDiamond}{Yen $!A$ iku !!a{formula}, mula $\Diamond !A$ iku
  !!a{formula}.}{}

\tagitem{limitClause}{Ora ana liyane kang kalebu !!a{formula}.}{}
\end{enumerate}
\end{defn}

\begin{tagblock}{defNot,defOr,defAnd,defIf,defIff,defTrue,defFalse,defBox,defDiamond}
\begin{defn}
Formula kang dibangun nganggo operator kang wis ditegesi kudu
dimangerteni kaya mangkene:

\begin{tagenumerate}{defTrue,defFalse,defNot,defOr,defAnd,defIf,defIff,defBox,defDiamond}

\tagitem{defTrue}{$\ltrue$ minangka cekakan saka
  \iftag{prvFalse}{$\lnot\lfalse$}{$(!A \lor \lnot !A)$ kanggo sawijining
    !!{formula} atomik tetep~$!A$}.}{}

\tagitem{defFalse}{$\lfalse$ minangka cekakan saka
  \iftag{prvTrue}{$\lnot\ltrue$}{$(!A \land \lnot !A)$ kanggo sawijining
    !!{formula} atomik tetep~$!A$}.}{}

\tagitem{defNot}{$\lnot !A$ minangka cekakan saka $!A \lif \lfalse$.}{}

\tagitem{defOr}{$!A \lor !B$ minangka cekakan saka
  \iftag{prvAnd}{$\lnot(\lnot !A \land \lnot !B)$}{$\lnot !A \lif
    !B$}.}{}

\tagitem{defAnd}{$!A \land !B$ minangka cekakan saka
  \iftag{prvOr}{$\lnot(\lnot !A \lor \lnot !B)$}{$\lnot (!A \lif
    \lnot !B)$}.}{}

\tagitem{defIf}{$!A \lif !B$ minangka cekakan saka
  \iftag{prvOr}{$\lnot !A \lor !B$}{$\lnot (!A \land \lnot !B)$}.}{}

\tagitem{defIff}{$!A \liff !B$ minangka cekakan saka $(!A \lif !B) \land (!B
  \lif !A)$.}{}

\tagitem{defBox}{$\Box !A$ minangka cekakan saka $\lnot\Diamond\lnot !A$. }{}

\tagitem{defDiamond}{$\Diamond !A$ minangka cekakan saka $\lnot\Box\lnot !A$.}{}
\end{tagenumerate}
\end{defn}
\end{tagblock}

Yen !!a{formula}~$!A$ ora ngemot $\Box$ utawa $\Diamond$, kita ngarani iki \emph{tanpa modal}.

\end{document}
